\documentclass[11pt]{article}

\usepackage[a4paper,margin=25mm]{geometry}
\usepackage[T1]{fontenc}
\usepackage[utf8]{inputenc}
\usepackage{lmodern}
\usepackage{microtype}
\usepackage{amsmath,amssymb,mathtools}
\usepackage{amsthm}
\usepackage{array}
\usepackage{booktabs}
\usepackage{enumitem}
\usepackage{tabularx}
\usepackage{xcolor}
\usepackage{hyperref}

\definecolor{artifactblue}{HTML}{1F4E79}
\definecolor{artifactgreen}{HTML}{1E6B45}
\definecolor{artifactred}{HTML}{8B2F2F}
\definecolor{artifactgray}{HTML}{F2F4F5}

\hypersetup{
  colorlinks=true,
  linkcolor=artifactblue,
  urlcolor=artifactblue,
  pdftitle={The Checked Sampler Prefix of the HAETAE Mode-2 Parent},
  pdfsubject={Actual-parent procedure transfer and certificate-conditioned sampler prefix}
}

\setlist[itemize]{leftmargin=1.5em,itemsep=0.25em,topsep=0.35em}
\setlist[enumerate]{leftmargin=1.7em,itemsep=0.25em,topsep=0.35em}
\renewcommand{\arraystretch}{1.18}
\setlength{\emergencystretch}{3em}

\theoremstyle{definition}
\newtheorem{definition}{Definition}
\theoremstyle{plain}
\newtheorem{theorem}{Theorem}
\theoremstyle{remark}
\newtheorem*{remark}{Remark}

\newcommand{\code}[1]{\nolinkurl{#1}}
\newcommand{\W}{\mathbb W_{64}}
\newcommand{\uint}[1]{\langle #1\rangle}
\newcommand{\Expand}{\mathsf{Expand}}
\newcommand{\PM}{\mathsf P_{\mathrm M}}
\newcommand{\PV}{\mathsf P_{\mathrm V}}
\newcommand{\PE}{\mathsf P_{\eta}}
\newcommand{\Prefix}{\mathsf{Prefix}_{2,3}}

\newenvironment{resultbox}
  {\begin{center}\begin{minipage}{0.94\linewidth}\color{artifactgreen}\hrule\vspace{0.6em}\color{black}}
  {\vspace{0.6em}\color{artifactgreen}\hrule\end{minipage}\end{center}}

\newenvironment{boundarybox}
  {\begin{center}\begin{minipage}{0.94\linewidth}\color{artifactred}\hrule\vspace{0.6em}\color{black}}
  {\vspace{0.6em}\color{artifactred}\hrule\end{minipage}\end{center}}

\title{The Checked Sampler Prefix of the HAETAE Mode-2 Parent\\
\large A mathematical reading of \code{TargetKeygenMode2Parent*.ec}}
\author{HAETAE reference EasyCrypt artifact}
\date{Artifact state checked 2026-07-23}

\begin{document}
\maketitle

\begin{abstract}
The actual extracted mode-2 key-generation target contains the complete
\code{_keypair_full_m23} procedure.  The new refinement first transfers the
previous seed-expansion, sampler-leaf, and caller results to procedures in
that actual parent module.  It then checks a proof-only program that follows
the parent's sampler prefix for parameters \(k=2\) and \(m=3\): one seed
expansion, six matrix polynomials, two uniform-vector polynomials, and the
five eta polynomials of the first retry attempt.  Deterministic finite-prefix
certificates name a sufficient termination witness for each of these thirteen
sampler leaves.  Under the combined certificate predicate, a pHoare theorem
proves probability-one termination of the whole proof-only prefix.  The
checked prefix is deliberately not the full keypair procedure and proves no
rejection-sampling distribution or key-generation security theorem.
\end{abstract}

\begin{resultbox}
\textbf{The result in one sentence.}
Previously checked finite-stream, range, frame, and certificate-conditioned
leaf results now apply to the corresponding procedures in the actual
mode-2-parent target, and a proof-only \(k=2,m=3\) prefix composes their
on-return semantics and certificate-conditioned probability-one termination
through the first five eta nonces.
\end{resultbox}

\section{The Actual Parent and the Checked Cut}

The deterministic mode-2 entry
\code{crypto_sign_keypair_internal_mode2_jazz} invokes
\code{_keypair_full_m23} with
\[
 (k,m,\textit{vkbytes},\textit{best\_count},\tau,\textit{rem},
   \textit{singular\_bound})
 = (2,3,992,5,58,24,611098).
\]
The generated EasyCrypt target preserves the following call path:
\[
\begin{array}{rcl}
\textit{seedbuf}&\leftarrow&
  \text{\code{_kp_expand_seedbuf}}(\textit{seedbuf},\textit{rawSeed}),\\
\textit{mat}&\leftarrow&
  \text{\code{_kp_polymatkm_expand_matA}}(\textit{mat},\textit{seedbuf},2,3),\\
\textit{avec}&\leftarrow&
  \text{\code{_kp_polyveck_expand_vecA}}(\textit{avec},\textit{seedbuf},2,3).
\end{array}
\]
It then sets \(\textit{counter}=0\) and enters a retry loop.  One attempt
samples
\[
\begin{array}{rcl}
s_1&\leftarrow&
  \text{\code{_kp_polyvec_expand_eta}}(s_1,\textit{seedbuf},
    \textit{counter},3),\\
s_2&\leftarrow&
  \text{\code{_kp_polyvec_expand_eta}}(s_2,\textit{seedbuf},
    \textit{counter}+3,2),
\end{array}
\]
and advances the counter by \(3+2=5\).  The full parent next calls
\code{_kp_m23_matrix}, \code{_keypair_finalize_m23}, and
\code{_singular_full}; a singular-value comparison decides whether another
attempt is needed.  On acceptance it selects bytes \(0\)--\(31\) and
\(96\)--\(127\) of the expanded seed and packs the public and secret keys.

The authored module \code{CheckedMode2ParentSamplerPrefix} reproduces only
the green cut below:
\[
\begin{array}{c}
\fcolorbox{artifactgreen}{artifactgray}{\parbox{0.82\textwidth}{\centering
seed expansion \(\rightarrow\) \(2\times3\) matrix expansion
\(\rightarrow\) two-slot uniform vector
\(\rightarrow\) first \(3+2\) eta slots}}\\[0.7em]
\textcolor{artifactred}{\Downarrow\quad\text{proof cut}}\\[0.7em]
\fcolorbox{artifactred}{artifactgray}{\parbox{0.82\textwidth}{\centering
matrix/NTT computation \(\rightarrow\) finalize
\(\rightarrow\) singular-value rejection and retries
\(\rightarrow\) packing}}
\end{array}
\]
Every call above the cut uses a procedure from
\code{KeygenMode2ParentTarget.M}.  Nevertheless, the authored prefix is a
separate proof observer; it is not extracted Jasmin.

\begin{boundarybox}
\textbf{Identity boundary.}
\code{CheckedMode2ParentSamplerPrefix.run} is not
\code{Parent._keypair_full_m23}.  A theorem about the former must not be cited
as semantics or termination of the latter.
\end{boundarybox}

\section{Seven Bridges into the Actual Parent Module}

Earlier refinements were stated about procedures in
\code{KeygenSamplerCallersTarget.M}.  The actual parent is extracted into the
separate module \code{KeygenMode2ParentTarget.M}.  Similar source text in two
targets is not, by itself, a logical bridge between these EasyCrypt modules.
The refinement therefore proves seven relational equivalences:

\begin{center}
\begin{tabularx}{\textwidth}{>{\raggedright\arraybackslash}p{0.42\textwidth}X}
\toprule
Actual-parent procedure & Bridge theorem\\
\midrule
\code{_kp_expand_seedbuf} &
\code{kp_expand_seedbuf_cross_equiv}\\
\code{_kp_poly_uniform_at_seedbuf_8192} &
\code{uniform8192_leaf_cross_equiv}\\
\code{_kp_poly_uniform_at_seedbuf_2048} &
\code{uniform2048_leaf_cross_equiv}\\
\code{_kp_poly_uniform_eta_at_seedbuf_2048} &
\code{eta2048_leaf_cross_equiv}\\
\code{_kp_polymatkm_expand_matA} &
\code{expand_matA_cross_equiv}\\
\code{_kp_polyveck_expand_vecA} &
\code{expand_vecA_cross_equiv}\\
\code{_kp_polyvec_expand_eta} &
\code{expand_eta_cross_equiv}\\
\bottomrule
\end{tabularx}
\end{center}

Each judgment starts the parent and sampler-target procedures with equal
arguments and proves equal results.  The proofs are syntactic procedure
simulations.  The dedicated gate separately verifies that the 24 generated
theories shared by the two extraction closures are byte-identical.

The equivalences transport the existing contracts to parent-qualified
theorems.  Four Hoare theorems cover the exact 128-byte seed expansion and its
uniform, eta, and key slices.  Three more cover the matrix, uniform-vector,
and eta-vector finite streams together with their range and aggregate frame
properties.  Finally, three pHoare theorems transport probability-one
termination of the large uniform leaf, small uniform leaf, and eta leaf under
their exact progress-certificate and capacity premises.

\begin{remark}
The transported caller Hoare theorems are partial-correctness judgments: they
describe a call if it returns.  The transported leaf pHoare theorems are
totality judgments, but only for one leaf call satisfying its explicit
certificate.  These are different claims.
\end{remark}

\section{The Mode-2 Certificate Bundle}

Let \(S\) be the raw 32-byte seed and let \(E=\Expand(S)\) be the exact
128-byte SHAKE256 expansion.  Three endpoint maps provide deterministic
finite witnesses:
\[
 L_{\mathrm M}:\{0,1\}\times\{0,1,2\}\to\mathbb Z,\qquad
 L_{\mathrm V}:\{0,1\}\to\mathbb Z,\qquad
 L_\eta:\{0,\ldots,4\}\to\mathbb Z.
\]
Write \(P_{\mathrm u}(E,o,n,L)\) for
\code{uniform_progress_prefix} and
\(P_\eta(E,o,n,L)\) for \code{eta_progress_prefix}.  The specification
packages the following finite obligations.

\begin{definition}[Matrix bundle]
\[
 \PM(E,L_{\mathrm M})\iff
 \bigwedge_{\substack{0\le i<2\\0\le j<3}}
 P_{\mathrm u}(E,0,256i+j,L_{\mathrm M}(i,j)).
\]
This is six certificates, one for each matrix polynomial.
\end{definition}

\begin{definition}[Uniform-vector bundle]
\[
 \PV(E,L_{\mathrm V})\iff
 \bigwedge_{0\le i<2}
 P_{\mathrm u}(E,0,515+i,L_{\mathrm V}(i)).
\]
The starting nonce \(515=256\cdot2+3\) is the caller schedule for
\(k=2,m=3\).
\end{definition}

\begin{definition}[First-attempt eta bundle]
\[
 \PE(E,L_\eta)\iff
 \bigwedge_{0\le i<5}P_\eta(E,32,i,L_\eta(i)).
\]
Slots \(0,1,2\) belong to \(s_1\), while slots \(3,4\) belong to \(s_2\).
\end{definition}

The predicate \code{mode2_sampler_prefix_progress} requires these three
bundles for every expanded buffer matching the raw seed.  Since
\code{output_matches} determines the exact SHAKE256 bytes, this form lets a
caller state the premise before executing seed expansion.

\begin{boundarybox}
\textbf{Certificate boundary.}
The specification defines what a suitable set of \(6+2+5=13\) finite
certificates means.  It does not prove that these witnesses exist for every
seed, and a deterministic certificate does not establish uniformity,
independence, or any other sampler-distribution property.
\end{boundarybox}

\section{On-Return Semantics of the Checked Prefix}

The procedure \code{CheckedMode2ParentSamplerPrefix.run} returns
\[
 (E,M,A,S_1,S_2,c).
\]
The Hoare theorem
\code{checked_mode2_parent_sampler_prefix_correct} proves, conditional on
return:
\begin{itemize}
\item \(E\) is the exact 128-byte expansion of the raw seed, including the
  named uniform, eta, and key slices;
\item \(M\) contains six exact finite SHAKE128-decoded polynomial witnesses,
  every produced coefficient is in the checked uniform range, and bytes
  outside the six-polynomial prefix retain their initial values;
\item \(A\) has the analogous finite-stream, range, and frame properties for
  two uniform polynomials;
\item \(S_1\) consists of three exact finite SHAKE256 eta-decoded polynomials,
  is centered, and satisfies its aggregate frame;
\item \(S_2\) has the same properties for two eta polynomials starting at
  nonce \(3\); and
\item \(c=5\), both as the mode-2 retry span and as a concrete 64-bit word.
\end{itemize}

This is a semantic composition result rather than merely a call-schedule
equivalence.  Its matrix, vector, and eta conclusions refer to the same
returned expanded seed \(E\), and its calls are made through the actual parent
module.

\section{Certificate-Conditioned Totality of the Prefix}

The totality proof is assembled in the same order as the program:
\begin{center}
\begin{tabularx}{\textwidth}{>{\raggedright\arraybackslash}p{0.48\textwidth}X}
\toprule
Program fragment & Probability-one termination theorem\\
\midrule
Exact parent seed expansion, with its output relation &
\code{parent_kp_expand_seedbuf_output_matches_pr}\\
Fixed \(2\times3\) matrix caller &
\code{mode2_matrix_uniform_progress_ll}\\
Fixed two-slot uniform-vector caller &
\code{mode2_expand_vecA_progress_ll}\\
Eta caller at nonce \(0\), count \(3\) &
\code{mode2_eta_nonce0_count3_progress_ll}\\
Eta caller at nonce \(3\), count \(2\) &
\code{mode2_eta_nonce3_count2_progress_ll}\\
Whole proof-only sampler prefix &
\code{checked_mode2_parent_sampler_prefix_progress_ll}\\
\bottomrule
\end{tabularx}
\end{center}

The matrix and vector judgments assume their respective certificate bundles.
The two eta judgments assume the shared five-slot first-attempt bundle.
The final pHoare theorem starts before seed expansion: its premise is exactly
\[
  \text{\code{mode2_sampler_prefix_progress}}
  (S,L_{\mathrm M},L_{\mathrm V},L_\eta).
\]
Seed expansion returns the exact \(E=\Expand(S)\); the specification then
specializes the premise to \(\PM(E,L_{\mathrm M})\),
\(\PV(E,L_{\mathrm V})\), and \(\PE(E,L_\eta)\).  The four fixed callers are
therefore lossless in sequence, so \code{run} returns with probability one.

The pHoare theorem's postcondition is \code{true}: it is a termination theorem,
not a duplicate semantic theorem.  Together with
\code{checked_mode2_parent_sampler_prefix_correct}, it gives the two separate
ingredients ``returns under the bundle'' and ``has the displayed result when
it returns.''  The artifact does not package them as one stronger pHoare
postcondition.

\section{Dependency and Proof Cuts}

The checked implications can be read from bottom to top:
\[
\begin{array}{c}
\fcolorbox{artifactgreen}{artifactgray}{\parbox{0.86\textwidth}{\centering
exact Keccak/SHAKE refinements; finite uniform and eta decoders;
leaf range/frame and certificate-conditioned leaf totality}}\\[0.65em]
\Downarrow\\[0.65em]
\fcolorbox{artifactgreen}{artifactgray}{\parbox{0.86\textwidth}{\centering
sampler-target seed and caller theorems}}\\[0.65em]
\Downarrow\quad\text{seven cross-module equivalences}\\[0.65em]
\fcolorbox{artifactgreen}{artifactgray}{\parbox{0.86\textwidth}{\centering
parent-qualified seed, caller, and leaf theorems}}\\[0.65em]
\Downarrow\quad\text{\(k=2,m=3\) proof-only composition}\\[0.65em]
\fcolorbox{artifactgreen}{artifactgray}{\parbox{0.86\textwidth}{\centering
checked first-attempt sampler-prefix semantics and
certificate-conditioned probability-one termination}}\\[0.65em]
\textcolor{artifactred}{\Downarrow}\quad
\text{\textcolor{artifactred}{matrix/NTT, finalize, singular rejection, retries, and packing not crossed}}\\[0.65em]
\fcolorbox{artifactred}{artifactgray}{\parbox{0.86\textwidth}{\centering
\code{Parent._keypair_full_m23} semantics and termination;
abstract key generation and implementation-level security}}
\end{array}
\]

The red cut is substantive.  In particular, the present certificate bundle
contains only eta nonces \(0,\ldots,4\).  A rejected first attempt would use
the next five nonces, and no theorem here proves acceptance, bounds the number
of retries, or supplies certificates for later attempts.

\section{Exact Claim Boundary}

\begin{center}
\begin{tabularx}{\textwidth}{>{\raggedright\arraybackslash}p{0.48\textwidth}X}
\toprule
Claim & Status\\
\midrule
The seven selected procedures in the actual parent target agree with their
sampler-target copies on equal inputs &
Proved by seven \code{equiv} judgments\\
Exact seed slices and finite-stream/range/frame caller semantics are available
for actual-parent-module procedures &
Proved as parent-qualified Hoare theorems\\
Each selected uniform or eta leaf terminates under its own exact progress
certificate and memory bounds &
Proved as parent-qualified pHoare theorems\\
The fixed matrix, uniform-vector, and two first-attempt eta callers terminate
under the corresponding bundle predicates &
Proved by the four fixed-caller pHoare theorems\\
The proof-only mode-2 prefix has the displayed semantic postcondition on
return &
Proved\\
The proof-only mode-2 prefix terminates under
\code{mode2_sampler_prefix_progress} &
Proved by \code{checked_mode2_parent_sampler_prefix_progress_ll}\\
Every seed has all thirteen progress certificates &
Not proved\\
The actual \code{_keypair_full_m23} returns, accepts its first attempt, or
satisfies the prefix postcondition &
Not proved\\
Matrix/NTT, finalize, singular-value rejection, later retry attempts, and key
packing are semantically refined &
Not proved\\
The decoded values have the distributions assumed by HAETAE, or the target
inherits the abstract EUF-CMA theorem &
Not proved\\
\bottomrule
\end{tabularx}
\end{center}

\section{Replaying the Machine-Checked Evidence}

From the \code{haetae-ref-easycrypt/} directory, run
\begin{center}
\code{./scripts/verify-keygen-mode2-parent-proof.sh}
\end{center}
The dedicated gate checks pinned source hashes, rejects extraction drift in
both the actual-parent and sampler-caller closures, verifies identity of their
24 shared generated theories, and compiles the 17-entry proof manifest with
\code{-no-eco}.  It also scans the authored specifications and refinements for
\code{admit}, \code{abort}, and axiom declarations.  The retained 2026-07-23
full run reports 17 of 17 entries compiled and all source, drift, shared-file
identity, proof-hole, and axiom checks passing.

Build this guide with
\begin{quote}
\begin{verbatim}
cd haetae-ref-easycrypt/latex
make build/TargetKeygenMode2Parent.pdf
\end{verbatim}
\end{quote}

\section*{Conclusion}

The refinement now reaches procedures in the actual mode-2-parent extraction
and composes one faithful sampler-prefix shape using the fixed dimensions and
first-attempt nonce schedule.  The combined certificate premise is sufficient
for probability-one termination of that proof-only prefix.  Its exact stopping
point is equally important: the checked program is proof-only, and the full
parent's algebraic work, singularity-controlled retry loop, packing,
distributions, and security composition remain outside the theorem.

\end{document}
